\documentclass[tikz,border=6mm]{standalone}

% ============================================================
%  Packages
% ============================================================
\usepackage[UTF8]{ctex}
\usepackage{xcolor}
\usepackage{tikz}
\usetikzlibrary{matrix,positioning,fit,backgrounds}

% ============================================================
%  Color Palette
% ============================================================
\definecolor{TitleBlue}{HTML}{1F4E79}
\definecolor{HeaderBlue}{HTML}{D9EAF7}
\definecolor{HeaderText}{HTML}{103A5C}
\definecolor{RowA}{HTML}{F8FBFE}
\definecolor{RowB}{HTML}{EEF6FC}
\definecolor{GridLine}{HTML}{B7CCDD}
\definecolor{AccentGreen}{HTML}{E8F5E9}
\definecolor{AccentOrange}{HTML}{FFF3E0}
\definecolor{AccentPurple}{HTML}{F3E5F5}
\definecolor{NoteBg}{HTML}{FAFCFF}
\definecolor{NoteBorder}{HTML}{B7CCDD}

% ============================================================
%  Layout Parameters
%  后续调整版面时，优先修改这里
% ============================================================
\newcommand{\TableGap}{8mm}
\newcommand{\TableRowHeight}{22mm}
\newcommand{\HeaderRowHeight}{13mm}
\newcommand{\TitleTextWidth}{16.7cm}
\newcommand{\NoteTextWidth}{16.7cm}
\newcommand{\NoteColWidth}{5.05cm}

% ============================================================
%  Typography Helpers
% ============================================================
\newcommand{\code}[1]{\texttt{#1}}

% ============================================================
%  Content: Title / Headers
% ============================================================
\newcommand{\TitleText}{AI 时代多语言开发模式快速参考表}
\newcommand{\SubTitleText}{AI 最适合做什么 / 人最该盯住什么 / 主要验证工具}

\newcommand{\HdrLang}{语言}
\newcommand{\HdrAI}{AI 最适合做}
\newcommand{\HdrHuman}{人最该盯住}
\newcommand{\HdrTools}{主要验证工具}

% ============================================================
%  Content: Table Rows
% ============================================================
\newcommand{\RowBashAI}{脚手架、批处理、文本处理、自动化}
\newcommand{\RowBashHuman}{quoting、错误传播、幂等性、破坏性命令}
\newcommand{\RowBashTools}{\code{shellcheck}、\code{bash -n}、测试环境}

\newcommand{\RowPyAI}{业务逻辑、自动化、测试、数据处理}
\newcommand{\RowPyHuman}{类型边界、异常路径、依赖、运行时行为}
\newcommand{\RowPyTools}{\code{pytest}、\code{mypy/pyright}、\code{ruff}}

\newcommand{\RowCAI}{驱动、底层接口、协议、资源控制}
\newcommand{\RowCHuman}{指针、生命周期、越界、UB、并发}
\newcommand{\RowCTools}{compiler warnings、ASan/UBSan、静态分析}

\newcommand{\RowCppAI}{复杂系统、抽象、并发、性能}
\newcommand{\RowCppHuman}{ownership、RAII、模板、线程安全、性能}
\newcommand{\RowCppTools}{clang-tidy、ASan/UBSan/TSan、benchmark}

\newcommand{\RowRustAI}{并发、系统编程、安全抽象}
\newcommand{\RowRustHuman}{\code{unsafe}、生命周期建模、API 设计、性能}
\newcommand{\RowRustTools}{\code{cargo clippy}、\code{cargo test}、Miri、sanitizer}

% ============================================================
%  TikZ Styles
% ============================================================
\tikzset{
    titlebox/.style={
        draw=TitleBlue,
        fill=TitleBlue!5,
        rounded corners=3mm,
        line width=0.9pt,
        inner xsep=6mm,
        inner ysep=4.5mm,
        text width=\TitleTextWidth,
        align=center
    },
    tablecell/.style={
        draw=GridLine,
        line width=0.6pt,
        align=center,
        anchor=center,
        inner xsep=4.5mm,
        inner ysep=4.5mm,
        font=\small
    },
    headercell/.style={
        fill=HeaderBlue,
        text=HeaderText,
        font=\bfseries\small,
        align=center,
        minimum height=\HeaderRowHeight
    },
    datacell/.style={
        align=center,
        minimum height=\TableRowHeight
    },
    notebox/.style={
        draw=NoteBorder,
        fill=NoteBg,
        rounded corners=2.5mm,
        line width=0.7pt,
        text width=\NoteTextWidth,
        align=left,
        inner xsep=6mm,
        inner ysep=4.5mm
    }
}

\begin{document}
\begin{tikzpicture}[font=\small]

% ============================================================
%  Title
% ============================================================
\node[titlebox] (titlebox) {
    {\Large\bfseries\color{TitleBlue}\TitleText}\\[1mm]
    {\normalsize\color{HeaderText}\SubTitleText}
};

% ============================================================
%  Main Table
%  注意：所有单元格均由 TikZ node 自身负责自动换行与居中，
%       不再使用顶对齐 parbox，因此垂直居中更稳定。
% ============================================================
\matrix[
    matrix of nodes,
    below=\TableGap of titlebox,
    nodes={tablecell},
    row sep=-\pgflinewidth,
    column sep=-\pgflinewidth,
    column 1/.style={
        nodes={text width=1.55cm,font=\bfseries\normalsize,align=center}
    },
    column 2/.style={nodes={text width=3.95cm,align=center}},
    column 3/.style={nodes={text width=4.25cm,align=center}},
    column 4/.style={nodes={text width=4.55cm,align=center}},
    row 1/.style={nodes={headercell}},
    row 2/.style={nodes={datacell,fill=RowA}},
    row 3/.style={nodes={datacell,fill=RowB}},
    row 4/.style={nodes={datacell,fill=RowA}},
    row 5/.style={nodes={datacell,fill=RowB}},
    row 6/.style={nodes={datacell,fill=RowA}}
] (tbl) {
    \HdrLang & \HdrAI & \HdrHuman & \HdrTools \\
    Bash   & \RowBashAI & \RowBashHuman & \RowBashTools \\
    Python & \RowPyAI   & \RowPyHuman   & \RowPyTools \\
    C      & \RowCAI    & \RowCHuman    & \RowCTools \\
    C++    & \RowCppAI  & \RowCppHuman  & \RowCppTools \\
    Rust   & \RowRustAI & \RowRustHuman & \RowRustTools \\
};

% ============================================================
%  Column Accent Overlays
% ============================================================
\begin{scope}[on background layer]
    \node[
        fit=(tbl-2-2)(tbl-6-2),
        fill=AccentGreen,
        fill opacity=0.16,
        rounded corners=1.5mm,
        inner sep=0pt,
        draw=none
    ] {};
    \node[
        fit=(tbl-2-3)(tbl-6-3),
        fill=AccentOrange,
        fill opacity=0.18,
        rounded corners=1.5mm,
        inner sep=0pt,
        draw=none
    ] {};
    \node[
        fit=(tbl-2-4)(tbl-6-4),
        fill=AccentPurple,
        fill opacity=0.18,
        rounded corners=1.5mm,
        inner sep=0pt,
        draw=none
    ] {};
\end{scope}

% ============================================================
%  Notes
% ============================================================
\node[notebox,below=\TableGap of tbl] (note) {
    {\bfseries\color{TitleBlue}使用提示}\\[1.5mm]
    \begin{minipage}[t]{\NoteColWidth}
        {\bfseries 1. 通用工作流}\\
        人定义目标与约束 $\rightarrow$ AI 生成/修改 $\rightarrow$ 工具验证 $\rightarrow$ 人 Review
    \end{minipage}
    \hfill
    \begin{minipage}[t]{\NoteColWidth}
        {\bfseries 2. 语言差异的本质}\\
        骨架相似，但每种语言的风险点、验证工具与 Review 重点不同。
    \end{minipage}
    \hfill
    \begin{minipage}[t]{\NoteColWidth}
        {\bfseries 3. 关键原则}\\
        让 AI 提高速度，让工具尽早暴露错误，让工程经验负责最终决策。
    \end{minipage}
};

\end{tikzpicture}
\end{document}
